% !TeX root = ../papers/nc-p3.tex
\section{二つのVeronese像の交わりと向きの分岐}
\label{sec:veronese-pencil-quotient}
$M$を十五次元成分の添字保持同型による有理商とする。$Q=A^2$の評価二十点を
$Y\subset\mathbb P(Q^*)$とし、それを含む二つのVeronese三次元多様体を
$V_A,V_C$と書く。以下の一般点に関する主張は標数零である。

\begin{theorem}\label{thm:veronese-pencil-quotient}
一般点では$V_A\cap V_C=Y$がスキームとして成立する。
奇位置の側でも同じであり、二つの二十点集合は互いの射影Gale変換である。
$M$から第一の四次鉛筆の射影同値類への関数体拡大は四元群をGalois群とする四次拡大である。
二鉛筆のうちどちらが第一かも忘れると、拡大の次数は8でGalois群は正方形の二面体群となる。
\end{theorem}
\begin{proof}
三次評価が同型なので$Y$を$V_A\simeq\mathbb P^3$で読むと、二十点の
イデアルの四次元数は$35-20=15$である。第一辺の六行四列の線型評価行列が期待余次元を持つため、
Eagon--Northcott分解によりイデアルはその四次最大小行列式で生成される。
$V_C$の二次小行列式を$V_A$へ制限すると、これらは二十点を消す四次式になる。
$t$の素数2657のファイバーでは、二十行十列の二次評価行列$X_A,X_C$について
$X_A T=\operatorname{diag}(s)X_C$を解く。$s$を決める百本の式は
$X_A$の左核を両辺に掛けた式であり、核は単元ベクトルの張る直線である。
得られた$T$を対称四行行列の全二次小行列式に入れ、四次の35単項式で展開すると
行の階数は15である。奇位置の側も15となる。
\path{computations/verify_configuration_intersections.py}が両方を検算する。
従って標数零の$t$および一般点でも全四次生成元を得る。これが交叉のスキーム構造まで決める。
一方$Q^\perp$の評価点は、$Q$の評価行列の関係を基底とした行列の二十個の対応行であり、
Gale変換の定義そのものである。各点の代表を変えても相手側の行倍率が変わるだけである。

一般鉛筆の中の行列式の集合は内在的である。既知の二本の共通三重点集合から上の
$Y$を回復できる。また鉛筆全体の共通三重点をまず取り、その任意の一点で
前節の偏極を課すこともできる。実際、二本の行列式では二階微分の成分が
対称四行行列の二次小行列式を張る。これは行列式の四重偏極で対角二方向または
非対角方向を選べば相補な二行二列式が得られることから従う（対称添字の
三通りの組合せの間には通常の直線関係がある）。従って集合としてはもちろん、
共通三重点のスキームも上記の交叉に一致する。

四位置の巡回移動を$T$、第一位置を固定した逆転を$J$とする。
逆転が$M$を保つことは前に示した。巡回移動は二つの隣接切断を行・列置換と
転置を除いて交換し、四元群の二つの薄い色を保つ。移動後に非退化チャートに
入ることは$t$で検算できる。実際、二つの切断の各厚い四行ブロックで四つの
組内・交差小行列式は全て非零で、薄い二行列の各成分も非零である。
薄い階数一行列の第0列を$u^2,u^3$とすると、規格化の$z$は
$u^2_2u^2_3u^3_0u^3_1/(u^3_2u^3_3u^2_0u^2_1)$であり、
移動後の第一切断で$-1/8$となる。列の選択や対角倍率はこの比に影響しない。
よって五次元の楕円曲線族の移動後の像は、非退化チャートの最大四次元の
例外表示には含まれず、同じ五次元族の閉包を支配する。$t$における全成分の一意性から
$T$は$M$を保つ。$T^4=J^2=1$, $JTJ=T^{-1}$であり、
四配置が一般に全て非同値であるため八作用は互いに異なる。

一つの鉛筆から二つの三重Veronese軌跡の一方を選ぶと、その上で交叉点を読むことにより
第一配置を射影同値を除いて復元できる。さらに向きを選べば前節により窓は決まる。
従って高々四通りであり、$1,T^2,J,T^2J$は全てその鉛筆を保つので丁度四通りである。
鉛筆は$Y$を決め、相手側はそのGale変換から回復できる。両鉛筆が射影同値なら、
対応するVeronese上で偶奇のいずれかの配置が射影同値になるので一般には相異なる。
従って順序なしの鉛筆対からのファイバーは八作用の一軌道で尽くされる。
標数零であるため、いずれの関数体拡大も表示した群の固定体拡大となる。
\end{proof}

\begin{corollary}\label{cor:orientation-ramification}
第一配置への二次拡大は、奇二位置が対称となる因子の一般点で分岐指数2を持つ。
その因子の一般配置における四辺有限のAS延長は一つである。
偶二位置が対称となる因子の一般点では逆転による二窓は異なる。
\end{corollary}
\begin{proof}
$r$を一位置移動して$r'=Tr$とする。前節の$r$の平方根スキームを
奇偶を交換して読むと、$r'$では第一側に二つの被約符号軌道があり、
第二側には一つの長さ2の軌道がある。従って第二空間は一意である。
第三位置で第一空間を再利用する候補は、前節の$r$における奇始点の三次交叉が
零である計算で排除される。

この計数は奇対称因子の一般点にも成立する。固有な平方根スキームのファイバーの
長さは$r'$の近傍で各側高々$2^{20}$である。第一側には各$2^{19}$点の相異なる
二つの軌道がある。第二側では二つの一般切断が因子上で一致するので、その根は
因子の一般点で不分岐ではない。これは二つの切断を番号と行平方根を選んで
エタール局所的に取れば、不分岐の場合の開対角に反するためである。従って
一つの符号軌道だけで既に全長さを使い切る。三次交叉零の条件は開であり、
実在する幾何学的一般AS窓だけが残る。

この因子は十四次元の片側対称族の巡回像であり、一般点では射影安定群が自明である。
四位置の点を番号付けし各位置の五点を枠に固定すれば、基底変換を除いた滑らかな
局所座標が得られる。番号の有限な選び直しは一般の安定群自明な点で自由なので、
商もそこで滑らかである。逆転$J$は因子を点ごとに固定し、全空間の一般点では
非自明である。因子に沿う座標の$J$平均と、横座標の反不変部分を取れる。
実際、滑らかな局所環上の有限位数作用は標数零で接空間の基底を固有ベクトルに
持ち上げられ、接空間上も恒等なら完備化上で恒等になってしまう。
従って横方向は符号表現で、局所不変環の横座標はその二乗となる。これは関数体の
因子付値についての指数を与える。射影像の正規化前の局所環を滑らかとは仮定していない。
偶対称因子の一般点では二奇配置の非同値性が保たれ、逆転は窓を固定しない。
\end{proof}
